% Options for packages loaded elsewhere
\PassOptionsToPackage{unicode}{hyperref}
\PassOptionsToPackage{hyphens}{url}
\documentclass[
  12pt,
]{article}
\usepackage{xcolor}
\usepackage{amsmath,amssymb}
\setcounter{secnumdepth}{-\maxdimen} % remove section numbering
\usepackage{iftex}
\ifPDFTeX
  \usepackage[T1]{fontenc}
  \usepackage[utf8]{inputenc}
  \usepackage{textcomp} % provide euro and other symbols
\else % if luatex or xetex
  \usepackage{unicode-math} % this also loads fontspec
  \defaultfontfeatures{Scale=MatchLowercase}
  \defaultfontfeatures[\rmfamily]{Ligatures=TeX,Scale=1}
\fi
\usepackage{lmodern}
\ifPDFTeX\else
  % xetex/luatex font selection
\fi
% Use upquote if available, for straight quotes in verbatim environments
\IfFileExists{upquote.sty}{\usepackage{upquote}}{}
\IfFileExists{microtype.sty}{% use microtype if available
  \usepackage[]{microtype}
  \UseMicrotypeSet[protrusion]{basicmath} % disable protrusion for tt fonts
}{}
\makeatletter
\@ifundefined{KOMAClassName}{% if non-KOMA class
  \IfFileExists{parskip.sty}{%
    \usepackage{parskip}
  }{% else
    \setlength{\parindent}{0pt}
    \setlength{\parskip}{6pt plus 2pt minus 1pt}}
}{% if KOMA class
  \KOMAoptions{parskip=half}}
\makeatother
\usepackage{color}
\usepackage{fancyvrb}
\newcommand{\VerbBar}{|}
\newcommand{\VERB}{\Verb[commandchars=\\\{\}]}
\DefineVerbatimEnvironment{Highlighting}{Verbatim}{commandchars=\\\{\}}
\newenvironment{Shaded}{}{}
\newcommand{\AlertTok}[1]{\textcolor[rgb]{1.00,0.00,0.00}{\textbf{#1}}}
\newcommand{\AnnotationTok}[1]{\textcolor[rgb]{0.38,0.63,0.69}{\textbf{\textit{#1}}}}
\newcommand{\AttributeTok}[1]{\textcolor[rgb]{0.49,0.56,0.16}{#1}}
\newcommand{\BaseNTok}[1]{\textcolor[rgb]{0.25,0.63,0.44}{#1}}
\newcommand{\BuiltInTok}[1]{\textcolor[rgb]{0.00,0.50,0.00}{#1}}
\newcommand{\CharTok}[1]{\textcolor[rgb]{0.25,0.44,0.63}{#1}}
\newcommand{\CommentTok}[1]{\textcolor[rgb]{0.38,0.63,0.69}{\textit{#1}}}
\newcommand{\CommentVarTok}[1]{\textcolor[rgb]{0.38,0.63,0.69}{\textbf{\textit{#1}}}}
\newcommand{\ConstantTok}[1]{\textcolor[rgb]{0.53,0.00,0.00}{#1}}
\newcommand{\ControlFlowTok}[1]{\textcolor[rgb]{0.00,0.44,0.13}{\textbf{#1}}}
\newcommand{\DataTypeTok}[1]{\textcolor[rgb]{0.56,0.13,0.00}{#1}}
\newcommand{\DecValTok}[1]{\textcolor[rgb]{0.25,0.63,0.44}{#1}}
\newcommand{\DocumentationTok}[1]{\textcolor[rgb]{0.73,0.13,0.13}{\textit{#1}}}
\newcommand{\ErrorTok}[1]{\textcolor[rgb]{1.00,0.00,0.00}{\textbf{#1}}}
\newcommand{\ExtensionTok}[1]{#1}
\newcommand{\FloatTok}[1]{\textcolor[rgb]{0.25,0.63,0.44}{#1}}
\newcommand{\FunctionTok}[1]{\textcolor[rgb]{0.02,0.16,0.49}{#1}}
\newcommand{\ImportTok}[1]{\textcolor[rgb]{0.00,0.50,0.00}{\textbf{#1}}}
\newcommand{\InformationTok}[1]{\textcolor[rgb]{0.38,0.63,0.69}{\textbf{\textit{#1}}}}
\newcommand{\KeywordTok}[1]{\textcolor[rgb]{0.00,0.44,0.13}{\textbf{#1}}}
\newcommand{\NormalTok}[1]{#1}
\newcommand{\OperatorTok}[1]{\textcolor[rgb]{0.40,0.40,0.40}{#1}}
\newcommand{\OtherTok}[1]{\textcolor[rgb]{0.00,0.44,0.13}{#1}}
\newcommand{\PreprocessorTok}[1]{\textcolor[rgb]{0.74,0.48,0.00}{#1}}
\newcommand{\RegionMarkerTok}[1]{#1}
\newcommand{\SpecialCharTok}[1]{\textcolor[rgb]{0.25,0.44,0.63}{#1}}
\newcommand{\SpecialStringTok}[1]{\textcolor[rgb]{0.73,0.40,0.53}{#1}}
\newcommand{\StringTok}[1]{\textcolor[rgb]{0.25,0.44,0.63}{#1}}
\newcommand{\VariableTok}[1]{\textcolor[rgb]{0.10,0.09,0.49}{#1}}
\newcommand{\VerbatimStringTok}[1]{\textcolor[rgb]{0.25,0.44,0.63}{#1}}
\newcommand{\WarningTok}[1]{\textcolor[rgb]{0.38,0.63,0.69}{\textbf{\textit{#1}}}}
\usepackage{longtable,booktabs,array}
\usepackage{calc} % for calculating minipage widths
% Correct order of tables after \paragraph or \subparagraph
\usepackage{etoolbox}
\makeatletter
\patchcmd\longtable{\par}{\if@noskipsec\mbox{}\fi\par}{}{}
\makeatother
% Allow footnotes in longtable head/foot
\IfFileExists{footnotehyper.sty}{\usepackage{footnotehyper}}{\usepackage{footnote}}
\makesavenoteenv{longtable}
\setlength{\emergencystretch}{3em} % prevent overfull lines
\providecommand{\tightlist}{%
  \setlength{\itemsep}{0pt}\setlength{\parskip}{0pt}}
\usepackage{bookmark}
\IfFileExists{xurl.sty}{\usepackage{xurl}}{} % add URL line breaks if available
\urlstyle{same}
\hypersetup{
  hidelinks,
  pdfcreator={LaTeX via pandoc}}

\author{SCX}
\date{2026}

\begin{document}

\section{Hostile Review Report: Spring Self-Evolution Convergence Analysis}\label{spring-hostile-review}

\begin{center}\rule{0.5\linewidth}{0.5pt}\end{center}

\textbf{Reviewer Positioning}: Rigorous non-convex optimization reviewer\\
\textbf{Review Date}: 2026-06-29\\
\textbf{Review Target}: \texttt{theory/self\_evolution/spring\_convergence\_analysis.md}\\
\textbf{Review Method}: Theorem-by-theorem attack + assumption completeness review + notation consistency review\\
\textbf{Review Standard}: Impolite. No acceptance of ``partially rigorous + conjecture.'' Every hole will be exploited.

\begin{center}\rule{0.5\linewidth}{0.5pt}\end{center}

\subsection{0. Overall Impression}\label{overall-impression}

This document claims to provide ``rigorous mathematical derivation + honest critique'' of Spring self-evolution dynamics convergence. After reading, the conclusion is: \textbf{The honesty labels are deceptive.} Multiple results labeled \texttt{[Rigorous Proof]} rely on undeclared assumptions or contain logical leaps; Theorem 1.4's constructive proof has an order-of-magnitude error; Theorem 2.1 ignores the core issue of variable action spaces; Theorem 3.1's counterexample uses hand-picked magic numbers rather than being derived from Spring dynamics. Below is a theorem-by-theorem dissection.

\begin{center}\rule{0.5\linewidth}{0.5pt}\end{center}

\subsection{1. Theorem 1.2's $O(\log T/\sqrt{T})$ Convergence Rate --- Full Attack on Ghadimi-Lan Conditions}\label{attack-thm-1.2}

\subsubsection{1.1 Applicability Review of the Ghadimi-Lan Framework}\label{ghadimi-lan-applicability}

Theorem 1.2 claims to use ``standard non-convex SGD analysis.'' The literature reference is Ghadimi \& Lan (2013), but the conditions under Spring's setting have not been carefully verified.

Ghadimi-Lan (2013, SIAM J. Optim.) core assumptions:
- \textbf{(GL1)} $f$ is $L$-smooth: $\|\nabla f(x) - \nabla f(y)\| \leq L\|x - y\|$ for all $x, y$.
- \textbf{(GL2)} Stochastic gradient $G(x, \xi)$ satisfies $\mathbb{E}[G(x,\xi)] = \nabla f(x)$ and $\mathbb{E}[\|G(x,\xi) - \nabla f(x)\|^2] \leq \sigma^2$.
- \textbf{(GL3)} $f$ is bounded below: $\inf_x f(x) > -\infty$.

\subsubsection{1.2 Attack Point \#1: Global $L$-Smoothness Does Not Hold in Spring (or the Constant Explodes)}\label{attack-1-global-smoothness}

Spring's loss landscape has $K!$ permutation symmetries (\S1.2). Consider two points $\theta$ and $\pi \cdot \theta$ in different permutation basins. Along the path between them, the loss function must cross ridges between basins. The height and curvature of these ridges determine the global Lipschitz constant.

\textbf{Question}: What is the global $L$-smoothness constant of the Spring loss function? If $L$ is related to $K!$ (e.g., $L = \Omega(K!)$ because it must cover ridges between all permutation basins), then:
- The condition $\alpha_t \leq 1/L$ requires step size $\leq 1/K!$, while the theorem uses $\alpha_t = \alpha/\sqrt{t}$, violating this condition for $t < L^2$.
- The constant factor $L\alpha^2\sigma^2/2$ in the convergence rate may render the bound \textbf{vacuous}.

\textbf{The document contains no estimate of the magnitude of $L$.} This is a serious omission --- without an estimate of $L$, the constant in $O(\log T/\sqrt{T})$ could be exponentially large.

\subsubsection{1.3 Attack Point \#2: The $\log T$ Factor Is Self-Inflicted --- Ghadimi-Lan Does Not Require It Originally}\label{attack-2-logT-self-inflicted}

Ghadimi-Lan (2013) standard result: using \textbf{constant step size} $\alpha_t = \alpha = \sqrt{\frac{2(f(x_1)-f^*)}{L\sigma^2 T}}$, one obtains:
\[\min_{t=1,\ldots,T} \mathbb{E}[\|\nabla f(x_t)\|^2] \leq \sqrt{\frac{2L\sigma^2(f(x_1)-f^*)}{T}} = O\left(\frac{1}{\sqrt{T}}\right)\]
\textbf{Note: No $\log T$ factor.}

The document's Theorem 1.2 has $\log T/\sqrt{T}$ because it uses decaying step size $\alpha_t = \alpha/\sqrt{t}$ --- $\sum \alpha_t^2 = \alpha^2 \sum 1/t = O(\log T)$ replaces $T\alpha^2$. But decaying step size provides \textbf{no benefit} in this analysis --- constant step size already yields $O(1/\sqrt{T})$ rate (no $\log T$) in non-convex SGD. The document's step size schedule \textbf{worsens} the convergence rate.

\textbf{Either the document is unaware of Ghadimi-Lan's constant step size result (cited but incorrectly used), or it deliberately chose a slower schedule to make the result appear more ``realistic'' --- both are unacceptable.}

\subsubsection{1.4 Attack Point \#3: Tension Between Bounded Variance Assumption and Permutation Symmetry}\label{attack-3-bounded-variance}

The assumption $\mathbb{E}[\|\zeta_t\|^2 | \mathcal{F}_t] \leq \sigma^2$ requires uniform upper bound on gradient noise. At the boundaries of permutation-symmetric basins, for samples from different basins, mini-batch gradients may point in completely different directions --- at which point the variance of $\zeta_t$ (difference between mini-batch and full-batch gradients) may \textbf{diverge at basin boundaries}.

This is not theoretical nitpicking --- it is a known practical phenomenon: when data contains multiple symmetric patterns, SGD's gradient noise is anomalously large at symmetry boundaries. Spring's $K!$ permutation symmetry amplifies this problem.

\textbf{The document acknowledges this} (\S1.2.2's ``honest note'' admits that noise covariance on the quotient space requires more careful analysis), but Theorem 1.2's statement and proof do not reflect this limitation. The theorem is labeled \texttt{[Rigorous Proof]} without noting that the implicit assumption of uniform noise variance may not hold.

\subsubsection{1.5 Attack Point \#4: Gradient Norm Convergence $\neq$ Loss Convergence}\label{attack-4-gradient-norm-vs-loss}

The document itself acknowledges this (``honest critique''), but the acknowledgment is insufficient. In highly non-convex landscapes, points with $\|\nabla f\| \approx 0$ include:
- Local minima (possibly shallow, high loss).
- Saddle points (loss may be arbitrarily high).
- Global minima (where we want to be).

For Spring, permutation symmetry implies $K!$ equivalent global minima --- but also exponentially many saddle points connecting them. Theorem 1.2's bound cannot distinguish ``converging to a global minimum'' from ``stuck at a saddle point.'' The document says ``this is a gradient norm bound and does not necessarily imply loss function value convergence'' --- this should be written in the theorem statement, not buried in a later ``honest critique.''

\subsubsection{1.6 Verdict on Theorem 1.2}\label{verdict-thm-1.2}

\begin{longtable}[]{@{}ll@{}}
\toprule\noalign{}
Issue & Severity \\
\midrule\noalign{}
\endhead
\bottomrule\noalign{}
\endlastfoot
Global $L$-smoothness in Spring is unestimated & \textbf{High} --- bound may be vacuous \\
$\log T$ factor is self-inflicted by step size choice & \textbf{Medium} --- constant step size avoids it \\
Bounded variance assumption is suspect at basin boundaries & \textbf{Medium} --- needs verification \\
Gradient norm convergence $\neq$ loss convergence is concealed & \textbf{Medium} --- insufficient honesty labeling \\
\end{longtable}

\textbf{Summary}: The \texttt{[Rigorous Proof]} label is technically correct (standard non-convex SGD analysis is indeed rigorous), but this is a rigorous proof for \textbf{some function that may not be Spring} --- whether Spring satisfies all conditions of the analysis has not been verified.

\begin{center}\rule{0.5\linewidth}{0.5pt}\end{center}

\subsection{2. Theorem 1.4's Information-Theoretic Lower Bound --- Line-by-Line Attack on the Constructive Example}\label{attack-thm-1.4}

\subsubsection{2.1 Attack Point \#1: $T^{1/4}$ Order-of-Magnitude Error in $\sigma^2/\sqrt{T}$}\label{attack-1-magnitude-error}

The document claims (\S1.4.1 proof):
\begin{quote}
``For $T$ steps, the noise-induced displacement is $\sigma \sqrt{\sum \alpha_t^2} \approx \sigma \cdot O(T^{1/4})$ (when $\alpha_t = 1/\sqrt{t}$).''
\end{quote}

\textbf{This is wrong.} $\sum_{t=1}^T \alpha_t^2 = \sum_{t=1}^T 1/t = H_T = O(\log T)$, so:
\[\sigma \sqrt{\sum_{t=1}^T \alpha_t^2} = \sigma \cdot O(\sqrt{\log T})\]
\textbf{not} $\sigma \cdot O(T^{1/4})$.

The origin of $O(T^{1/4})$ is baffling. If $\alpha_t = 1/\sqrt{t}$, then $\alpha_t^2 = 1/t$, sum is $O(\log T)$, square root is $O(\sqrt{\log T})$. To get $O(T^{1/4})$, one would need $\alpha_t = 1/t^{3/4}$ or a similar step size --- but the theorem statement explicitly writes $\alpha_t = \alpha/\sqrt{t}$.

\textbf{This order-of-magnitude error appears in the core computation of the theorem's proof.} The entire noise-driven displacement argument depends on this magnitude, hence the derivation of the $\sigma^2/\sqrt{T}$ term in the lower bound is unreliable.

\subsubsection{2.2 Attack Point \#2: 1D Construction Cannot Represent the High-Dimensional Spring Landscape}\label{attack-2-1d-construction}

The document constructs $f(x) = \sin^2(\pi K x / 2)$ on $[0,1]$ as a 1D function, claiming it ``simulates $K$ equivalent stationary points of $K$ heads under permutation.''

\textbf{Questions}:
- The 1D function has $K/2$ global minima ($f=0$), uniformly spaced. Spring's high-dimensional landscape has $K!$ permutation-symmetric stationary points. The scaling of $K/2$ and $K!$ is completely different --- for $K=4$, the construction has 2 minima, but Spring has 24. Why does the lower bound not reflect the $K!$ factor?
- A 1D function cannot simulate the \textbf{dimensionality} and \textbf{connectivity structure} between basins in a high-dimensional landscape. In high dimensions, SGD can escape between basins via saddle points (which is exactly what Proposition 1.3 shows!), but such escape is impossible in a 1D function (must climb over barriers).
- Projecting a high-dimensional problem onto a 1D ``order parameter'' implies an undeclared averaging over gradient noise and gradient signal. This averaging may systematically alter the convergence rate.

\textbf{The number of minima in the 1D construction ($K/2$) and in Spring ($K!$) are not at the same order of magnitude.} Hence the $K$-dependence of the lower bound is unverified.

\subsubsection{2.3 Attack Point \#3: Is the Function Family Within Spring's Hypothesis Class?}\label{attack-3-function-family}

$f(x) = \sin^2(\pi K x / 2)$ is:
- Smooth ($C^\infty$) \checkmark
- Periodic $\times$ (Spring's loss landscape is not periodic --- although it has permutation symmetry, the distribution of different permutation basins in parameter space is not a uniform lattice but a complex geometry described by Weyl chambers).
- Bounded below \checkmark
- All minima equivalent (all global) \checkmark (consistent with permutation symmetry).

But Spring's loss function is not a simple scaling of $\sin^2$ --- it has $K!$ minima distributed in an $O(K d_s^2)$-dimensional space, and the connectivity structure between minima involves a complex saddle network. To use a 1D periodic function to ``lower-bound'' such a complex landscape, the argument requires a \textbf{dimensionality reduction lemma} to formalize it --- which the document does not provide.

\subsubsection{2.4 Attack Point \#4: Suspicious $\min$ Structure of the Lower Bound $\Omega(\min(\sigma^2/\sqrt{T}, 1/T))$}\label{attack-4-min-structure}

For high noise ($\sigma^2 \gg 1/\sqrt{T}$), the lower bound is $\Omega(\sigma^2/\sqrt{T})$. For low noise, it becomes $\Omega(1/T)$. This means in the noiseless limit ($\sigma \to 0$), the lower bound is $\Omega(1/T)$ --- which is the rate of \textbf{first-order deterministic gradient descent}.

But deterministic gradient descent on non-convex smooth functions can achieve $O(1/T)$ gradient norm convergence. So the lower bound in the low-noise case is $\Omega(1/T)$, matching the upper bound --- superficially tight. But the $\min$ operation here is essentially saying ``take the smaller of the two as the lower bound,'' and the structure $\min(\sigma^2/\sqrt{T}, 1/T)$ means: \textbf{when $T$ is not large enough, the noise term dominates; when $T$ is large enough, the deterministic error dominates.} The correctness of this structure depends critically on the magnitude analysis in \S2.1 --- and we have already found an error there.

\subsubsection{2.5 Verdict on Theorem 1.4}\label{verdict-thm-1.4}

\begin{longtable}[]{@{}
  >{\raggedright\arraybackslash}p{(\linewidth - 2\tabcolsep) * \real{0.4286}}
  >{\raggedright\arraybackslash}p{(\linewidth - 2\tabcolsep) * \real{0.5714}}@{}}
\toprule\noalign{}
\begin{minipage}[b]{\linewidth}\raggedright
Issue
\end{minipage} & \begin{minipage}[b]{\linewidth}\raggedright
Severity
\end{minipage} \\
\midrule\noalign{}
\endhead
\bottomrule\noalign{}
\endlastfoot
$\sigma\sqrt{\sum\alpha_t^2} = O(T^{1/4})$ is an order-of-magnitude error & \textbf{Fatal} --- should be $O(\sqrt{\log T})$ \\
1D construction: $K/2$ minima vs Spring's $K!$ mismatch & \textbf{High} --- $K$-scaling unverified \\
Projection from $K!$-dim to 1D is unargued & \textbf{High} --- dimensionality reduction unproven \\
Function family (periodic $\sin^2$) mismatches Spring's non-periodic basin geometry & \textbf{Medium} --- needs stronger analogy argument \\
\end{longtable}

\textbf{Summary}: The $T^{1/4}$ order-of-magnitude error alone suffices to invalidate Theorem 1.4's proof. Even if corrected, the chasm between the 1D construction and Spring's high-dimensional landscape remains unbridged. The \texttt{[Rigorous Proof]} label in its current form is misleading.

\begin{center}\rule{0.5\linewidth}{0.5pt}\end{center}

\subsection{3. Theorem 2.1 Regret Upper Bound --- Full Review of Exp3 Algorithm Adaptation}\label{attack-thm-2.1}

\subsubsection{3.1 Attack Point \#1: Action Space Variability --- Hidden Cracks in Standard Exp3 Analysis}\label{attack-1-action-space}

Exp3's standard regret bound (Auer et al., 2002) assumes a \textbf{fixed action space} ($N$ arms, always available). Spring's memory bank dynamics introduce two problems:

\textbf{(a) Variable action space size}: $|\mathcal{M}_t^{\text{dormant}}|$ varies with $t$. Memory bank pruning removes state atoms; discovery of new states adds atoms. Exp3's weight vector $w_t$ changes dimension accordingly. The document implicitly replaces $|\mathcal{M}_t^{\text{dormant}}|$ with $|\mathcal{S}|$ (total number of state atoms) --- but this only gives an upper bound when $|\mathcal{M}_t^{\text{dormant}}| \leq |\mathcal{S}|$. This is a valid relaxation, but the relaxation magnitude is $|\mathcal{S}| / |\mathcal{M}_t^{\text{dormant}}|$; when most state atoms are ``active'' (non-dormant), the bound becomes very loose.

\textbf{(b) Weight initialization for new arms}: When new state atoms are discovered and added to $\mathcal{M}_{t+1}^{\text{dormant}}$, Exp3 must initialize a weight for them. The document does not specify the initialization rule. If the initial weight is too small (relative to existing arms), the new arm won't be selected for a long time --- increasing regret. If the initial weight is too large, it may crowd out the selection probability of existing arms.

\subsubsection{3.2 Attack Point \#2: Adversarial Assumption --- Is the Adversary Oblivious?}\label{attack-2-adversarial}

The document claims memory bank dynamics produce \textbf{adversarial} rewards (\S2.2.1): ``Since memory bank dynamics depend on historical choices, the reward sequence may be adversarial (rather than stochastic i.i.d.).''

Exp3's regret bound indeed applies to \textbf{non-stochastic} (adversarial) reward sequences. But there is a subtlety:

In Auer et al.~(2002)'s framework, the adversary can be \textbf{adaptive} (decide current reward based on past choices), but cannot decide reward based on \textbf{the current round's randomization}. Spring's memory bank dynamics fall into the former category --- after selecting $a_t$ in round $t$, the memory bank update $\mathcal{M}_{t+1}$ affects round $t+1$'s reward but not round $t$'s (since round $t$'s reward is determined by the quality of round $t$'s state atoms, fixed after selecting $a_t$).

\textbf{But there is a circular causality problem}: Round $t$'s selection $a_t$ is based on some evaluation result, and this evaluation result \textbf{itself depends on the current Spring parameters $\theta_t$}. $\theta_t$ is determined by the memory bank states of the previous $t-1$ rounds through SE-1 (SGD update). Therefore the reward sequence depends on history through a complex feedback loop --- this is indeed an adversarial setting that Exp3 can handle. \checkmark

\textbf{A more subtle problem}: Exp3 requires rewards $r_t(a_t) \in [0,1]$. The document says rewards are ``quality scores from Yajie audit (normalized Cercis score or estimated detection margin $\Delta_s$).'' But estimating $\Delta_s$ requires \textbf{true labels} to determine $p_{\text{clean}}$ and $p_{\text{noisy}}$. In a real Spring system, Yajie audit only provides $C(s) = \sum v_m$ --- the expert vote consistency count, not the true detection margin. Replacing $\Delta_s$ with $\hat{\Delta}_s$ introduces estimation error, which is not modeled in the regret bound.

\subsubsection{3.3 Attack Point \#3: Comparator $\mu^*$ Definition Problem}\label{attack-3-comparator}

\[\mu^* = \max_{\mathcal{I} \subset \mathcal{S}, |\mathcal{I}| = B} \frac{1}{B} \sum_{s \in \mathcal{I}} \Delta_s\]

This comparator assumes:
1. \textbf{The optimal fixed set} is time-invariant --- but $\Delta_s$ itself varies over time ($\Delta_s(t)$), because Spring's parameters are evolving. The document uses $\Delta_s$ (no time index), implying a static $\Delta_s$ --- contradicting the entire time-dependent analysis of $\Delta_s(t)$ in \S3.
2. \textbf{All $\Delta_s$ are known in advance} --- the optimal fixed policy needs to know the true values of all $\Delta_s$. These values are unknown even in hindsight (because they depend on counterfactual Spring evolution trajectories). So $\mu^*$ is an \textbf{uncomputable} comparator.

In practice, the only available comparator is ``empirical estimate of the ex-post optimum,'' not $\mu^*$.

\subsubsection{3.4 Attack Point \#4: Undefined Behavior of Exp3's Importance Weighting upon Memory Bank Deletion}\label{attack-4-importance-weighting}

Exp3's update rule (\S2.2.1):
\[w_{t+1}(s) = w_t(s) \cdot \exp\left(\frac{\gamma \cdot \hat{r}_t(s)}{|\mathcal{M}_t^{\text{dormant}}| \cdot p_t(s)}\right)\]

When state $s$ is \textbf{pruned and removed} from the memory bank, what should happen to $w_{t+1}(s)$? If deleted, that arm can never be selected again --- but this means the pruning operation is \textbf{irreversible}, and if a high-quality state is mistakenly deleted (false negative), the algorithm has no opportunity to recover.

The document's Corollary 2.1 indeed discusses conservative pruning to avoid mistaken deletion --- but this turns pruning from a part of the algorithm into a constraint that must satisfy a specific inequality. If these constraints are not met (hard to guarantee in practice), the Exp3 regret bound needs an additional term to penalize mistaken deletions.

\subsubsection{3.5 Verdict on Theorem 2.1}\label{verdict-thm-2.1}

\begin{longtable}[]{@{}
  >{\raggedright\arraybackslash}p{(\linewidth - 2\tabcolsep) * \real{0.4286}}
  >{\raggedright\arraybackslash}p{(\linewidth - 2\tabcolsep) * \real{0.5714}}@{}}
\toprule\noalign{}
\begin{minipage}[b]{\linewidth}\raggedright
Issue
\end{minipage} & \begin{minipage}[b]{\linewidth}\raggedright
Severity
\end{minipage} \\
\midrule\noalign{}
\endhead
\bottomrule\noalign{}
\endlastfoot
Variable action space size --- Exp3 analysis assumes fixed & \textbf{Medium} --- Upper-bounding with $|\mathcal{S}|$ works but causes looseness \\
Reward signal depends on unobservable $\Delta_s$ & \textbf{High} --- In practice, the reward is a noisy estimate \\
Comparator $\mu^*$ uses static $\Delta_s$, contradicts \S3 & \textbf{Medium} --- Notation inconsistency \\
Weight management for added/deleted arms is undefined & \textbf{Medium} --- Needs algorithmic refinement \\
\end{longtable}

\textbf{Summary}: The \texttt{[Rigorous Proof]} label is correct for standard Exp3 analysis, but Spring's specificities (variable action space, estimation error, irreversible pruning) introduce complexities not covered by the standard analysis. Theorem 2.1 should be labeled \texttt{[Rigorous Proof (under assumptions of fixed action space and exact reward observation)]}.

\begin{center}\rule{0.5\linewidth}{0.5pt}\end{center}

\subsection{4. Theorem 3.1's Monotonicity Counterexample --- Legality Attack on the Construction}\label{attack-thm-3.1}

\subsubsection{4.1 Attack Point \#1: The Probability Numbers in the Counterexample Are Hand-Picked Magic Numbers}\label{attack-1-magic-numbers}

Theorem 3.1's construction gives specific numbers:
- $p_{\text{clean},a}(t) = 0.2 \to p_{\text{clean},a}(t+1) = 0.35$
- $p_{\text{noisy},a}(t) = 0.7 \to p_{\text{noisy},a}(t+1) = 0.65$

Where do these numbers come \textbf{from}? The document claims the causal chain is: ``Spring's self-attention mechanism updates parameters $\theta_t$ at round $t$ to minimize self-reconstruction loss\ldots Due to multi-head attention's shared representation, overfitting to noisy $s_a$ instances affects the representation of all instances of state $s_a$.''

But \textbf{there is no derivation} connecting the gradient descent of the self-reconstruction loss to specific numerical changes in expert disagreement probabilities. The statement that $p_{\text{clean}}$ changes from 0.2 to 0.35 is merely an \textbf{assertion}, without:
- Gradient computation of the Spring loss.
- Expert model response function to representation changes.
- Mapping from parameter changes to disagreement probability changes.

\subsubsection{4.2 Attack Point \#2: The Counterexample's Core Mechanism --- ``Overfitting to Noisy Samples'' --- Is Inconsistent with Spring's Training Objective}\label{attack-2-overfitting-mechanism}

\textbf{Spring's training objective is to minimize self-reconstruction error}:
\[\mathcal{L}_{\text{Spring}}(\theta) = \sum_{i=1}^{N} \|\text{Spring}_{\text{MH}}(s_i; \theta) - s_i\|^2\]
This is an \textbf{unsupervised} autoencoding objective. Spring is trained to reconstruct state atoms from context --- it \textbf{does not directly} use Cercis scores or label information.

Theorem 3.1's counterexample claims Spring ``overfits to mislabeled noisy samples in the memory bank.'' But ``mislabeling'' is the result of Cercis scoring (or Yajie audit) --- not an input to Spring training. Spring can only see state atoms $s_i$, not whether they are clean or noisy. How then does Spring ``overfit to noisy samples''?

\textbf{Verdict}: The counterexample conflates two distinct mechanisms --- the gatekeeper's quality assessment and Spring's self-reconstruction --- without establishing a causal link between them. The numerical assertion has no mathematical derivation to back it. The \texttt{[Counterexample Construction]} label overstates the level of rigor.

\begin{center}\rule{0.5\linewidth}{0.5pt}\end{center}

\subsection{5. Summary Verdict}\label{summary-verdict}

\begin{longtable}[]{@{}lll@{}}
\toprule\noalign{}
Theorem & Claimed Rigor & Actual Assessment \\
\midrule\noalign{}
\endhead
\bottomrule\noalign{}
\endlastfoot
Thm 1.1 (PL convergence) & Rigorous (under PL) & Conditionally correct; PL assumption unverified \\
Thm 1.2 (General non-convex) & Rigorous & Correct analysis, but may be vacuous for Spring (no $L$ estimate) \\
Thm 1.3 (Distance lower bound) & Rigorous & Correct \\
Thm 1.4 (Info-theoretic lower bound) & Rigorous & \textbf{Fatal order-of-magnitude error}; proof invalid \\
Thm 2.1 (Regret upper bound) & Partially rigorous & Correct for Exp3, but Spring's specificities not modeled \\
Thm 3.1 (Monotonicity counterexample) & Counterexample & Hand-picked numbers; no derivation from Spring dynamics \\
\end{longtable}

\textbf{Overall}: The document makes a valiant attempt at rigorous analysis of a complex system, but the honesty labels overstate the level of rigor. The $T^{1/4}$ error in Theorem 1.4 is particularly concerning and must be corrected. The most valuable contribution is the structural analysis of permutation symmetry in \S1.2.1--1.2.2, which is genuinely rigorous and novel.

\end{document}
